\documentclass[journal]{IEEEtran}
\usepackage[T1]{fontenc}
\usepackage{amsmath,amssymb,amsfonts}
\usepackage{graphicx}
\usepackage{booktabs}
\usepackage{multirow}
\usepackage{siunitx}
\usepackage{xcolor}
\usepackage{url}
\usepackage[hidelinks]{hyperref}
\graphicspath{{../../docs/figures/}}
\newcommand{\method}{EG-SWM}
% Result numbers come from scripts/eval/aggregate.py (tables/numbers.tex), as in the main paper; a missing key prints ??
\newcommand{\val}[1]{\ifcsname val@#1\endcsname\csname val@#1\endcsname\else\textcolor{red}{\textbf{??}}\fi}
\InputIfFileExists{tables/numbers}{}{}
\renewcommand{\thesection}{S\arabic{section}}
\renewcommand{\thetable}{S\arabic{table}}
\renewcommand{\thefigure}{S\arabic{figure}}

\begin{document}
\title{Supplementary Material for\\ Embodiment-Grounded Surgical World Models}
\author{Ajo~Babu~George}
\markboth{IEEE Transactions on Medical Robotics and Bionics, supplementary material}{}
\maketitle

\section{Pilot experiments}\label{sec:pilots}
The design of \method{}'s objective and the scene-action variant were chosen in pilots of size S, trained for 12k steps on 213,985 windows of a subset of the training parts (the first pilot of Table~\ref{tab:pilot-objective} on an earlier subset of 181,585 windows) and diagnosed on the same 256 validation windows. \emph{Replay} generates the 7 future latents from 2 context latents (10 Euler steps, the same noise for every model) with the actions inferred from the real clip and reports the mean squared error to the real latents. The pilots use another budget and other validation windows than the full runs, so their numbers are not comparable with those of the main paper.

\textbf{Objective.} Trained on the diffusion loss alone with $\beta=10^{-3}$, all eight dimensions of each arm action became inactive (Table~\ref{tab:pilot-objective}): replacing every arm's action with another window's changed the replay error by 0.001\,\%, while replacing the explicit endoscope action raised the loss by up to 14\,\% across noise levels. With $\beta=10^{-4}$, six dimensions stay active but the diffusion model hardly uses them. Adding the forward decoder (same data and $\beta$) lowers the replay error by 12\,\% (17\,\% inside the instrument regions) and makes the arm actions matter three times as much, mainly inside the regions where they are written. The forward decoder, not the weaker bottleneck, therefore makes the arm actions useful; the full-scale ablation without it is in the main paper.

\begin{table}[h]
\centering
\caption{Pilots: the objective of the learned actions (step 12k, the same 256 validation windows). Increase: rise of the replay error when the learned or the endoscope action of another window is used.}
\label{tab:pilot-objective}
\resizebox{\columnwidth}{!}{%
\begin{tabular}{lcccc}
\toprule
 & \multicolumn{2}{c}{Arms, diffusion loss only} & Arms & Monolithic \\
 & $\beta=10^{-3}$ & $\beta=10^{-4}$ & + decoder & + decoder \\
\midrule
Active arm-action dimensions & 0/8 & 6/8 & 8/8 & -- \\
KL per arm (nats) & 0.004 & 1.41 & 14.2 & 27.3$^\ast$ \\
Replay error, all positions & 0.609 & 0.605 & 0.534 & 0.445 \\
\quad inside instrument regions & 0.734 & 0.702 & 0.584 & 0.560 \\
\quad outside & 0.587 & 0.588 & 0.526 & 0.425 \\
Increase, learned actions (\%) & 0.001 & 3.1 & 10.7 & 24.3 \\
\quad inside / outside (\%) & 0.0 / 0.0 & 4.0 / 3.0 & 20.5 / 8.8 & 23.3 / 24.5 \\
Increase, endoscope action (\%) & 7.1 & 7.2 & 8.2 & 3.1 \\
\bottomrule
\multicolumn{5}{l}{\footnotesize $^\ast$ per transition, for one 24-D action}
\end{tabular}}
\end{table}

\textbf{Frame-wide learned actions and the endoscope.} The monolithic action replays 17\,\% more accurately than the arm actions, almost all of it outside the instrument regions, and takes over part of the explicit endoscope action (Table~\ref{tab:pilot-scene}). A fifth, scene action, pooled and written only outside every arm's region, closes the gap while the arm actions stay local, but it too weakens the endoscope action, even with every learned action zeroed, the condition under which endoscope control is evaluated. Dropping the learned actions alone during training (probability 0.2) restored the endoscope action only in part and cost replay accuracy. Inferring every learned action from camera-stabilised latents, $z_{t-1}$ warped onto $z_t$ by the measured endoscope motion, did not help, and the scene action then carried as much endoscope information as without stabilisation, so it does not simply copy the motion it sees. Neither variant was run at full scale.

\begin{table}[h]
\centering
\caption{Pilots: frame-wide learned actions and the endoscope action (step 12k, the same 256 validation windows). Swap: rise of the replay error with another window's action; for the endoscope with the learned actions zeroed.}
\label{tab:pilot-scene}
\resizebox{\columnwidth}{!}{%
\begin{tabular}{lccccc}
\toprule
 & Arms & Monolithic & Arms & + learned-only & + stabilised \\
 & + decoder & + decoder & + scene & dropout & inference \\
\midrule
Replay error & 0.534 & 0.445 & 0.444 & 0.479 & 0.458 \\
Endoscope swap (\%) & 7.4 & 3.0 & 3.2 & 4.7 & 3.4 \\
Arm swap, inside / outside (\%) & 20.5 / 8.8 & -- & 14.9 / 5.2 & 12.8 / 4.5 & -- \\
Endoscope motion from the scene action, $R^2$ & -- & -- & 0.16 / 0.22 / 0.04 & -- & 0.17 / 0.22 / 0.03 \\
\bottomrule
\multicolumn{6}{l}{\footnotesize $R^2$ of ridge regression to the measured $\Delta x$ / $\Delta y$ / $\log s$}
\end{tabular}}
\end{table}

\section{SurgWM re-implementation}\label{sec:surgwm}
We re-implement SurgWM from its paper and, for the latent action model, from the authors' released code. A VQ-VAE with spatio-temporal transformer encoder and decoder turns each frame at $120\times180$ pixels into $30\times45$ tokens from 1,024 codes. The latent action model works at $40\times60$ pixels and is causal over time: a learned prompt token of each frame is quantised to one of twelve actions, from which and the earlier frames its decoder rebuilds the next frame. The MaskGIT dynamics (12 blocks of width 512; rotary position embeddings instead of a learned one, so 52\,M parameters instead of 62.5\,M) predicts the tokens of a frame from those of up to 15 earlier frames and the action into the frame, in 25 decoding steps at temperature 1. As released, both quantisers move their codes by exponential moving averages and restart every code that falls out of use from an encoder output of the batch. Each model is trained for 60k steps at batch size 8 with the paper's learning rates ($10^{-4}$ for the dynamics, for which it gives none) on the first 16 frames of our training windows at 10\,fps; the token cross-entropy of the dynamics stops falling after about 10k steps.

On the test windows the tokenizer reconstructs the frames at \val{main/surgwm/recon_psnr}\,dB PSNR, \val{main/surgwm/recon_lpips} LPIPS and FVD \val{main/surgwm/recon_fvd} (the VAE of every other model: \val{main/recon/psnr}\,dB, \val{main/recon/lpips} and \val{main/recon/fvd}), so its resolution sets the LPIPS and FVD of its predictions more than their PSNR. Its latent action model has collapsed. Its \val{use/surgwm/n_actions} codes lie within \val{use/surgwm/code_spread}\,\% of their norm of each other (the tokenizer's codes: median \val{use/surgwm/tok_spread}\,\%); over the \val{use/surgwm/n_windows} test clips it infers one action for \val{use/surgwm/top_share}\,\% of the transitions and never uses \val{use/surgwm/unused} of the \val{use/surgwm/n_actions} (perplexity \val{use/surgwm/perplexity}); and whichever code its decoder is given, its error on each of \val{use/surgwm/dyn_windows} of these clips changes by at most \val{use/surgwm/code_change}\,\%, although zero vectors, far from every code, change it by \val{use/surgwm/zero_change}\,\% on average. The perplexity over training batches did not show the collapse, because restarted codes are near copies of the one in use. Given the earlier frames, the decoder rebuilds the next frame \val{use/surgwm/lam_below_copy}\,\% more accurately than copying the previous one, whatever the action: at 10\,fps the next frame leaves an action little to explain. The dynamics accordingly ignores the actions: for a fully masked frame after the real frames before it, another clip's actions raise its token cross-entropy by \val{use/surgwm/ce_swap}\,\% (\val{use/surgwm/dyn_windows} test clips), and replaying the inferred actions is worth \val{main/surgwm/delta_psnr}\,dB PSNR over another clip's. At SurgWM's own 1\,fps an action would have more to carry, but the model would no longer predict the 2.8\,s that every other model predicts; SurgWM's twelve discrete actions trained within our model, interface (iv) of the main paper, are used and controllable, and are the comparison of the action interface.

\section{Sensitivity}\label{sec:sensitivity}
\textbf{Sampler.} The main paper samples with 50 Euler steps without guidance. On 200 of the test windows, spread over the \val{main/n_cases} test cases, 100 steps change PSNR by at most \val{sens/s100_max}\,dB for any interface, whereas 20 steps cost \method{} \val{sens/S_arms/s20_loss}\,dB (\val{sens/S_arms/end20_loss}\,dB at 2.8\,s, Table~\ref{tab:sampler}). Guidance on the actions (scale 2) adds \val{sens/g2_range}\,dB for every interface and strengthens action following (\method{}'s $\Delta$PSNR \val{sens/S_arms/delta_g2} instead of \val{sens/S_arms/delta_g1}\,dB), but it is left out so that all interfaces are sampled alike. Neither choice changes the order of the interfaces.

\IfFileExists{tables/e1_sampler.tex}{%
\begin{table}[h]
\centering
\caption{Sampler: Euler steps and guidance scale on the actions (200 test windows).}
\label{tab:sampler}
\resizebox{\columnwidth}{!}{\begin{tabular}{lccccccc}
\toprule
Model & Euler steps & Guidance & PSNR $\uparrow$ & SSIM $\uparrow$ & LPIPS $\downarrow$ & PSNR at 2.8 s $\uparrow$ & $\Delta$PSNR $\uparrow$ \\
\midrule
EG-SWM & 20 & 1 & 21.98 & 0.699 & 0.163 & 19.49 & 2.94 \\
EG-SWM & 50 & 1 & 22.27 & 0.702 & 0.154 & 20.26 & 2.94 \\
EG-SWM & 100 & 1 & 22.28 & 0.702 & 0.152 & 20.32 & 2.94 \\
EG-SWM & 50 & 2 & 22.33 & 0.707 & 0.150 & 20.25 & 3.49 \\
Monolithic latent + camera & 20 & 1 & 22.58 & 0.714 & 0.148 & 20.51 & 3.14 \\
Monolithic latent + camera & 50 & 1 & 22.59 & 0.712 & 0.146 & 20.69 & 3.14 \\
Monolithic latent + camera & 100 & 1 & 22.56 & 0.711 & 0.146 & 20.70 & 3.13 \\
Monolithic latent + camera & 50 & 2 & 22.66 & 0.718 & 0.142 & 20.75 & 3.66 \\
Slots + camera (FLAM-style) & 20 & 1 & 22.98 & 0.725 & 0.142 & 20.94 & 3.69 \\
Slots + camera (FLAM-style) & 50 & 1 & 22.93 & 0.722 & 0.141 & 21.03 & 3.65 \\
Slots + camera (FLAM-style) & 100 & 1 & 22.90 & 0.721 & 0.141 & 21.03 & 3.63 \\
Slots + camera (FLAM-style) & 50 & 2 & 22.98 & 0.727 & 0.137 & 21.07 & 4.20 \\
Monolithic VQ + camera (SurgWM-style) & 20 & 1 & 21.03 & 0.673 & 0.187 & 18.54 & 2.09 \\
Monolithic VQ + camera (SurgWM-style) & 50 & 1 & 21.47 & 0.681 & 0.174 & 19.52 & 2.19 \\
Monolithic VQ + camera (SurgWM-style) & 100 & 1 & 21.49 & 0.681 & 0.173 & 19.60 & 2.20 \\
Monolithic VQ + camera (SurgWM-style) & 50 & 2 & 21.65 & 0.686 & 0.170 & 19.73 & 2.76 \\
EG-SWM + scene action & 20 & 1 & 22.43 & 0.710 & 0.152 & 20.23 & 3.08 \\
EG-SWM + scene action & 50 & 1 & 22.51 & 0.709 & 0.148 & 20.58 & 3.05 \\
EG-SWM + scene action & 100 & 1 & 22.49 & 0.709 & 0.148 & 20.60 & 3.04 \\
EG-SWM + scene action & 50 & 2 & 22.61 & 0.715 & 0.143 & 20.61 & 3.64 \\
\bottomrule
\end{tabular}
}
\end{table}}{}

\textbf{Context length.} With one or three context latents instead of two, on the same 200 windows, the frames that all three settings predict (0.9--3.2\,s) gain \val{ctx/per_latent}\,dB PSNR per added latent for every interface; each added latent also brings the last context frame 0.4\,s nearer. The order of the interfaces stays the same, and \method{} trails the monolithic latent by \val{ctx/S_global-S_arms/1}, \val{ctx/S_global-S_arms/2} and \val{ctx/S_global-S_arms/3}\,dB with one, two and three latents (Table~\ref{tab:context}).

\IfFileExists{tables/e1_context.tex}{%
\begin{table}[h]
\centering
\caption{Context length: PSNR over the frames that every setting predicts (0.9--3.2\,s; 200 test windows).}
\label{tab:context}
\resizebox{\columnwidth}{!}{\begin{tabular}{lccc}
\toprule
Model & 1 latent (0.1 s) & 2 latents (0.5 s) & 3 latents (0.9 s) \\
\midrule
EG-SWM & 20.30 & 21.56 & 22.87 \\
Monolithic latent + camera & 20.73 & 21.91 & 23.20 \\
Slots + camera (FLAM-style) & 21.05 & 22.25 & 23.52 \\
Monolithic VQ + camera (SurgWM-style) & 19.45 & 20.75 & 22.08 \\
EG-SWM + scene action & 20.53 & 21.81 & 23.16 \\
\bottomrule
\end{tabular}
}
\end{table}}{}

\textbf{Command size.} Endoscope commands in the main paper are the 99th percentile of the real frame-to-frame motion. On 100 of the control windows, commands of half, one and two times this size are realised at \val{scale/S_arms/0.5}, \val{scale/S_arms/1} and \val{scale/S_arms/2} of their magnitude (mean over the six commands), in the commanded direction in at least \val{scale/S_arms/sign_min}\,\% of the windows for every command and size (Table~\ref{tab:command}). Twice the 99th percentile lies outside almost all real endoscope motion, and the models conditioned on robot state or on twelve discrete latent actions fall short there alike (\val{scale/S_state/2} and \val{scale/S_global_vq/2}). Zooms are realised at \val{scale/S_arms/zoom_range}\,\% of their magnitude at every size, pans and tilts at \val{scale/S_arms/pantilt_range}\,\%.

\IfFileExists{tables/e1_camera_scale.tex}{%
\begin{table}[h]
\centering
\caption{Command size: mean endoscope gain over the six commands and the share of windows moving in the commanded direction (100 control windows).}
\label{tab:command}
\resizebox{0.8\columnwidth}{!}{\begin{tabular}{lccc}
\toprule
Model & Command scale & Mean gain & Right direction \\
\midrule
EG-SWM & $\times$0.5 & 0.79 & 99.7\,\% \\
EG-SWM & $\times$1 & 0.71 & 99.8\,\% \\
EG-SWM & $\times$2 & 0.57 & 100.0\,\% \\
Monolithic VQ + camera (SurgWM-style) & $\times$0.5 & 0.86 & 99.7\,\% \\
Monolithic VQ + camera (SurgWM-style) & $\times$1 & 0.82 & 99.7\,\% \\
Monolithic VQ + camera (SurgWM-style) & $\times$2 & 0.67 & 99.5\,\% \\
Robot state + camera & $\times$0.5 & 0.87 & 100.0\,\% \\
Robot state + camera & $\times$1 & 0.80 & 100.0\,\% \\
Robot state + camera & $\times$2 & 0.65 & 100.0\,\% \\
\bottomrule
\end{tabular}
}
\end{table}}{}

\section{Additional results}\label{sec:more}
The tables below are generated with those of the main paper and appear as their experiments finish: whether a learned action describes what the moving endoscope reveals; fidelity by the number of instruments bound to an arm; intervals over test cases for locality and for endoscope control and embodiment consistency, paired against \method{}; the paired differences quoted in the text; 14\,s autoregressive rollouts; \method{} with fewer instrument detections at test time; run-to-run variation over seeds; and the cost of training and prediction.

\IfFileExists{tables/e1_scene_reveal.tex}{%
\begin{table}[h]
\centering
\caption{Does a learned action describe what the moving endoscope reveals? Share of the change caused by replacing the action that falls in the grid cells the endoscope motion reveals, over their share of the grid (median over the \val{reveal/n} test windows in which they cover at least \val{reveal/min_area}\,\% of it; smaller and larger endoscope motion separately). Near 1: the effect ignores the revealed content.}
\resizebox{\columnwidth}{!}{\begin{tabular}{lcccc}
\toprule
Action replaced & Windows & Revealed-cell ratio & smaller motion & larger motion \\
\midrule
EG-SWM + scene: scene action & 121 & 1.03 & 1.05 & 1.03 \\
EG-SWM + scene: arm actions & 121 & 0.86 & 0.81 & 0.90 \\
EG-SWM: arm actions & 121 & 0.91 & 0.89 & 0.96 \\
Monolithic latent + camera & 121 & 1.01 & 1.00 & 1.02 \\
Slots + camera & 121 & 0.99 & 0.98 & 1.02 \\
\bottomrule
\end{tabular}
}
\end{table}}{}

\IfFileExists{tables/e1_binding.tex}{%
\begin{table}[h]
\centering
\caption{PSNR on the \val{main/n_windows} test windows grouped by how many mounted instruments are bound to an arm region for most of the predicted steps.}
\resizebox{\columnwidth}{!}{\begin{tabular}{lccc}
\toprule
 & \multicolumn{3}{c}{Instruments bound (windows)} \\
\cmidrule(lr){2-4}
Model & none (93) & one (179) & $\geq$2 (728) \\
\midrule
Copy last context frame & 23.79 & 20.68 & 21.54 \\
None (unconditional) & 20.27 & 18.08 & 19.01 \\
Robot state & 20.63 & 18.28 & 19.37 \\
Robot state + camera & 22.36 & 19.53 & 20.42 \\
Monolithic latent (Genie) & 24.97 & 21.59 & 22.39 \\
Monolithic latent + camera & 25.10 & 21.58 & 22.35 \\
Monolithic VQ + camera (SurgWM-style) & 23.73 & 20.36 & 21.23 \\
Slots + camera (FLAM-style) & 25.32 & 21.85 & 22.67 \\
EG-SWM, no camera action & 20.11 & 20.42 & 21.96 \\
EG-SWM & 22.94 & 20.83 & 22.17 \\
EG-SWM + scene action & 24.56 & 21.27 & 22.37 \\
\bottomrule
\end{tabular}
}
\end{table}}{}

\IfFileExists{tables/e1_locality.tex}{%
\begin{table}[h]
\centering
\caption{Locality of the learned actions with 95\,\% intervals over test cases, and paired against \method{} of size S (cases in which \method{}'s arm actions are the more local).}
\resizebox{\columnwidth}{!}{\begin{tabular}{lccc}
\toprule
Model & Locality$\uparrow$ & $\Delta$ vs EG-SWM & EG-SWM more local \\
\midrule
Monolithic latent (Genie) & 1.43\,{\scriptsize[1.32, 1.51]} & $-$0.99\,{\scriptsize[$-$1.25, $-$0.88]} & 19/19 \\
Monolithic latent + camera & 1.46\,{\scriptsize[1.31, 1.56]} & $-$0.96\,{\scriptsize[$-$1.22, $-$0.83]} & 19/19 \\
SurgWM (re-implementation) & 1.15\,{\scriptsize[1.10, 1.22]} & $-$1.27\,{\scriptsize[$-$1.56, $-$1.10]} & 19/19 \\
Monolithic VQ + camera (SurgWM-style) & 1.53\,{\scriptsize[1.41, 1.62]} & $-$0.89\,{\scriptsize[$-$1.15, $-$0.77]} & 19/19 \\
Slots + camera (FLAM-style) & 1.44\,{\scriptsize[1.37, 1.51]} & $-$0.98\,{\scriptsize[$-$1.23, $-$0.86]} & 19/19 \\
EG-SWM, no camera action & 2.02\,{\scriptsize[1.82, 2.17]} & $-$0.40\,{\scriptsize[$-$0.62, $-$0.31]} & 17/19 \\
EG-SWM & 2.42\,{\scriptsize[2.29, 2.70]} & -- & -- \\
EG-SWM + scene action & 3.50\,{\scriptsize[3.06, 3.82]} & 1.09\,{\scriptsize[0.71, 1.25]} & 0/19 \\
EG-SWM, diffusion loss only & 2.21\,{\scriptsize[2.09, 2.40]} & $-$0.21\,{\scriptsize[$-$0.42, $-$0.10]} & 10/19 \\
EG-SWM, arm actions not grounded & 1.84\,{\scriptsize[1.71, 2.03]} & $-$0.58\,{\scriptsize[$-$0.76, $-$0.48]} & 19/19 \\
EG-SWM, wrong arm binding & 0.90\,{\scriptsize[0.80, 0.99]} & $-$1.52\,{\scriptsize[$-$1.87, $-$1.34]} & 19/19 \\
EG-SWM + console readout & 2.40\,{\scriptsize[2.22, 2.63]} & $-$0.01\,{\scriptsize[$-$0.16, 0.07]} & 12/19 \\
EG-SWM, boxes held across detector gaps & 2.47\,{\scriptsize[2.31, 2.74]} & 0.05\,{\scriptsize[$-$0.06, 0.15]} & 10/19 \\
EG-SWM, size B & 2.49\,{\scriptsize[2.33, 2.77]} & 0.07\,{\scriptsize[$-$0.09, 0.17]} & 5/19 \\
Monolithic latent + camera, size B & 1.47\,{\scriptsize[1.42, 1.58]} & $-$0.95\,{\scriptsize[$-$1.20, $-$0.83]} & 19/19 \\
\bottomrule
\end{tabular}
}
\end{table}}{}

\IfFileExists{tables/e1_control.tex}{%
\begin{table*}[t]
\centering
\caption{Endoscope gain and embodiment consistency with 95\,\% intervals over test cases; $\Delta$: \method{} of size S minus the model, paired over the same windows (cases in which \method{} is better).}
\resizebox{\textwidth}{!}{\begin{tabular}{lcccccc}
\toprule
Model & Cam. gain $\uparrow$ & ECS-P $\uparrow$ & ECS-R $\uparrow$ & Halluc.\,\% $\downarrow$ & $\Delta$gain & $\Delta$halluc. \\
\midrule
EG-SWM & 0.70\,{\scriptsize[0.69, 0.71]} & 0.94\,{\scriptsize[0.92, 0.95]} & 0.81\,{\scriptsize[0.78, 0.84]} & 7.6\,{\scriptsize[4.9, 11.0]} & -- & -- \\
EG-SWM, diffusion loss only & 0.69\,{\scriptsize[0.68, 0.70]} & 0.94\,{\scriptsize[0.92, 0.95]} & 0.80\,{\scriptsize[0.76, 0.83]} & 7.4\,{\scriptsize[4.4, 11.1]} & 0.01\,{\scriptsize[0.00, 0.02]} (17/19) & 0.3\,{\scriptsize[$-$0.6, 1.1]} (7/19) \\
EG-SWM, wrong arm binding & 0.70\,{\scriptsize[0.70, 0.71]} & 0.94\,{\scriptsize[0.92, 0.96]} & 0.81\,{\scriptsize[0.79, 0.84]} & 7.6\,{\scriptsize[4.5, 11.6]} & $-$0.01\,{\scriptsize[$-$0.01, $-$0.00]} (5/19) & 0.0\,{\scriptsize[$-$0.8, 0.7]} (7/19) \\
EG-SWM + console readout & 0.70\,{\scriptsize[0.69, 0.71]} & 0.94\,{\scriptsize[0.92, 0.95]} & 0.82\,{\scriptsize[0.79, 0.84]} & 7.4\,{\scriptsize[4.5, 11.0]} & 0.00\,{\scriptsize[$-$0.01, 0.01]} (13/19) & 0.3\,{\scriptsize[$-$0.4, 1.0]} (10/19) \\
EG-SWM, boxes held across detector gaps & 0.65\,{\scriptsize[0.65, 0.66]} & 0.93\,{\scriptsize[0.91, 0.95]} & 0.82\,{\scriptsize[0.79, 0.84]} & 8.1\,{\scriptsize[5.1, 11.6]} & 0.04\,{\scriptsize[0.04, 0.05]} (19/19) & $-$0.4\,{\scriptsize[$-$1.0, 0.2]} (14/19) \\
EG-SWM, size B & 0.78\,{\scriptsize[0.77, 0.78]} & 0.95\,{\scriptsize[0.93, 0.96]} & 0.87\,{\scriptsize[0.84, 0.89]} & 6.5\,{\scriptsize[3.7, 10.2]} & $-$0.08\,{\scriptsize[$-$0.09, $-$0.07]} (0/19) & 1.2\,{\scriptsize[0.2, 2.1]} (6/19) \\
Monolithic latent + camera, size B & 0.18\,{\scriptsize[0.17, 0.18]} & 0.95\,{\scriptsize[0.93, 0.97]} & 0.87\,{\scriptsize[0.84, 0.90]} & 6.1\,{\scriptsize[3.1, 10.4]} & 0.52\,{\scriptsize[0.51, 0.53]} (19/19) & 1.5\,{\scriptsize[0.3, 2.5]} (3/19) \\
\bottomrule
\end{tabular}
}
\end{table*}}{}

\IfFileExists{tables/e1_contrasts.tex}{%
\begin{table*}[t]
\centering
\caption{Paired differences on the E1 windows ($a-b$, means over test cases with 95\,\% intervals, and cases in which $a$ is better).}
\resizebox{\textwidth}{!}{\begin{tabular}{lcccccc}
\toprule
$a - b$ & $\Delta$PSNR & $a$ better & $\Delta$SSIM & $a$ better & $\Delta$LPIPS & $a$ better \\
\midrule
EG-SWM $-$ Copy last context frame & 0.41\,{\scriptsize[0.24, 0.56]} & 17 & 0.010\,{\scriptsize[0.005, 0.015]} & 14 & 0.008\,{\scriptsize[0.006, 0.010]} & 1 \\
EG-SWM $-$ None (unconditional) & 3.04\,{\scriptsize[2.86, 3.23]} & 19 & 0.108\,{\scriptsize[0.103, 0.114]} & 19 & $-$0.085\,{\scriptsize[$-$0.090, $-$0.081]} & 19 \\
Slots + camera (FLAM-style) $-$ EG-SWM & 0.76\,{\scriptsize[0.66, 0.87]} & 19 & 0.022\,{\scriptsize[0.019, 0.024]} & 19 & $-$0.015\,{\scriptsize[$-$0.018, $-$0.013]} & 19 \\
Monolithic latent + camera $-$ EG-SWM & 0.46\,{\scriptsize[0.37, 0.57]} & 19 & 0.012\,{\scriptsize[0.010, 0.014]} & 19 & $-$0.010\,{\scriptsize[$-$0.012, $-$0.007]} & 19 \\
Monolithic latent (Genie) $-$ EG-SWM & 0.49\,{\scriptsize[0.39, 0.59]} & 19 & 0.013\,{\scriptsize[0.011, 0.015]} & 19 & $-$0.010\,{\scriptsize[$-$0.012, $-$0.008]} & 19 \\
EG-SWM + scene action $-$ EG-SWM & 0.37\,{\scriptsize[0.30, 0.45]} & 19 & 0.010\,{\scriptsize[0.008, 0.011]} & 19 & $-$0.008\,{\scriptsize[$-$0.009, $-$0.006]} & 19 \\
EG-SWM $-$ Monolithic VQ + camera (SurgWM-style) & 0.70\,{\scriptsize[0.62, 0.79]} & 19 & 0.020\,{\scriptsize[0.018, 0.022]} & 19 & $-$0.019\,{\scriptsize[$-$0.021, $-$0.017]} & 19 \\
EG-SWM $-$ Robot state + camera & 1.57\,{\scriptsize[1.45, 1.69]} & 19 & 0.039\,{\scriptsize[0.036, 0.042]} & 19 & $-$0.042\,{\scriptsize[$-$0.044, $-$0.039]} & 19 \\
Monolithic VQ + camera (SurgWM-style) $-$ Copy last context frame & $-$0.29\,{\scriptsize[$-$0.46, $-$0.12]} & 4 & $-$0.010\,{\scriptsize[$-$0.015, $-$0.005]} & 2 & 0.027\,{\scriptsize[0.026, 0.029]} & 0 \\
EG-SWM $-$ EG-SWM, no camera action & 0.49\,{\scriptsize[0.39, 0.61]} & 19 & 0.019\,{\scriptsize[0.015, 0.023]} & 19 & $-$0.012\,{\scriptsize[$-$0.015, $-$0.009]} & 19 \\
EG-SWM $-$ EG-SWM, diffusion loss only & 0.36\,{\scriptsize[0.32, 0.39]} & 19 & 0.011\,{\scriptsize[0.010, 0.012]} & 19 & $-$0.008\,{\scriptsize[$-$0.009, $-$0.007]} & 19 \\
EG-SWM, no camera action $-$ Robot state & 2.22\,{\scriptsize[2.01, 2.44]} & 19 & 0.078\,{\scriptsize[0.071, 0.085]} & 19 & $-$0.060\,{\scriptsize[$-$0.065, $-$0.056]} & 19 \\
Robot state $-$ None (unconditional) & 0.33\,{\scriptsize[0.27, 0.39]} & 19 & 0.011\,{\scriptsize[0.009, 0.013]} & 19 & $-$0.012\,{\scriptsize[$-$0.015, $-$0.010]} & 19 \\
Robot state + camera $-$ Robot state & 1.15\,{\scriptsize[1.03, 1.26]} & 19 & 0.058\,{\scriptsize[0.054, 0.061]} & 19 & $-$0.031\,{\scriptsize[$-$0.034, $-$0.029]} & 19 \\
Monolithic latent (Genie) $-$ Monolithic latent + camera & 0.02\,{\scriptsize[$-$0.00, 0.05]} & 11 & 0.001\,{\scriptsize[0.000, 0.002]} & 16 & $-$0.001\,{\scriptsize[$-$0.001, $-$0.000]} & 14 \\
EG-SWM, diffusion loss only $-$ EG-SWM & $-$0.36\,{\scriptsize[$-$0.39, $-$0.32]} & 0 & $-$0.011\,{\scriptsize[$-$0.012, $-$0.010]} & 0 & 0.008\,{\scriptsize[0.007, 0.009]} & 0 \\
EG-SWM, arm actions not grounded $-$ EG-SWM & 0.06\,{\scriptsize[0.03, 0.09]} & 14 & 0.001\,{\scriptsize[0.000, 0.002]} & 10 & $-$0.001\,{\scriptsize[$-$0.002, 0.000]} & 10 \\
EG-SWM, wrong arm binding $-$ EG-SWM & $-$0.02\,{\scriptsize[$-$0.04, $-$0.00]} & 7 & $-$0.001\,{\scriptsize[$-$0.001, 0.000]} & 5 & 0.001\,{\scriptsize[0.000, 0.002]} & 2 \\
EG-SWM + console readout $-$ EG-SWM & $-$0.02\,{\scriptsize[$-$0.05, $-$0.00]} & 8 & $-$0.000\,{\scriptsize[$-$0.001, 0.000]} & 8 & $-$0.000\,{\scriptsize[$-$0.001, 0.001]} & 9 \\
EG-SWM, boxes held across detector gaps $-$ EG-SWM & 0.05\,{\scriptsize[0.04, 0.07]} & 18 & 0.001\,{\scriptsize[0.001, 0.002]} & 17 & $-$0.001\,{\scriptsize[$-$0.002, $-$0.001]} & 16 \\
EG-SWM, size B $-$ EG-SWM & 0.43\,{\scriptsize[0.38, 0.47]} & 19 & 0.016\,{\scriptsize[0.015, 0.017]} & 19 & $-$0.012\,{\scriptsize[$-$0.013, $-$0.011]} & 19 \\
EG-SWM, size B $-$ Copy last context frame & 0.83\,{\scriptsize[0.66, 1.00]} & 17 & 0.026\,{\scriptsize[0.021, 0.031]} & 19 & $-$0.004\,{\scriptsize[$-$0.006, $-$0.001]} & 13 \\
Monolithic latent + camera $-$ EG-SWM, size B & 0.04\,{\scriptsize[$-$0.07, 0.15]} & 9 & $-$0.004\,{\scriptsize[$-$0.006, $-$0.002]} & 4 & 0.002\,{\scriptsize[0.000, 0.005]} & 5 \\
Slots + camera (FLAM-style) $-$ EG-SWM, size B & 0.34\,{\scriptsize[0.23, 0.45]} & 17 & 0.005\,{\scriptsize[0.003, 0.008]} & 16 & $-$0.003\,{\scriptsize[$-$0.005, $-$0.001]} & 14 \\
EG-SWM, size B $-$ Monolithic VQ + camera (SurgWM-style) & 1.13\,{\scriptsize[1.02, 1.23]} & 19 & 0.036\,{\scriptsize[0.034, 0.038]} & 19 & $-$0.031\,{\scriptsize[$-$0.033, $-$0.029]} & 19 \\
\bottomrule
\end{tabular}
}
\end{table*}}{}

\IfFileExists{tables/e1_long.tex}{%
\begin{table*}[t]
\centering
\caption{14\,s autoregressive rollouts: PSNR per 2.8\,s chunk and generated instruments that the robot log does not mount, early and late.}
\resizebox{\textwidth}{!}{\begin{tabular}{lcccccccc}
\toprule
Model & PSNR 0.0--2.8 s & PSNR 2.8--5.6 s & PSNR 5.6--8.4 s & PSNR 8.4--11.2 s & PSNR 11.2--14.0 s & LPIPS 11.2--14.0 s & Unmounted at 0.6 s & Unmounted at 13.6 s \\
\midrule
Copy last context frame & 22.34 & 20.02 & 18.92 & 18.27 & 18.13 & 0.286 & -- & -- \\
EG-SWM & 22.56 & 19.38 & 18.17 & 17.64 & 17.20 & 0.352 & 0.099 & 0.152 \\
Monolithic latent + camera & 22.95 & 19.78 & 18.57 & 17.97 & 17.56 & 0.342 & 0.120 & 0.217 \\
Slots + camera (FLAM-style) & 23.29 & 20.11 & 18.89 & 18.26 & 17.84 & 0.333 & 0.098 & 0.184 \\
Monolithic VQ + camera (SurgWM-style) & 21.81 & 18.60 & 17.63 & 17.07 & 16.72 & 0.380 & 0.084 & 0.180 \\
Robot state + camera & 20.41 & 17.04 & 16.15 & 15.82 & 15.56 & 0.427 & 0.075 & 0.267 \\
None (unconditional) & 18.76 & 15.44 & 14.68 & 14.42 & 14.26 & 0.505 & 0.084 & 0.657 \\
EG-SWM + scene action & 22.94 & 19.62 & 18.49 & 17.88 & 17.44 & 0.349 & 0.121 & 0.187 \\
\bottomrule
\end{tabular}
}
\end{table*}}{}

\IfFileExists{tables/e2_probes.tex}{%
\begin{table*}[t]
\centering
\caption{Kinematic grounding, every target: $R^2$ of linear probes from frozen features to the end-effector translation (pos), rotation (rot) and jaw change of the same arm over each latent step, on held-out episodes (mean over the two instrument arms).}
\resizebox{\textwidth}{!}{\begin{tabular}{llccccccccc}
\toprule
 & & \multicolumn{3}{c}{SutureBot} & \multicolumn{3}{c}{CRCD} & \multicolumn{3}{c}{JIGSAWS} \\
\cmidrule(lr){3-5}\cmidrule(lr){6-8}\cmidrule(lr){9-11}
Features & Setting & pos & rot & jaw & pos & rot & jaw & pos & rot & jaw \\
\midrule
Latent change in the arm region & model-free & 0.40 & 0.07 & 0.07 & 0.00 & $-$0.00 & $-$0.01 & 0.05 & $-$0.00 & 0.00 \\
Latent change over the frame & model-free & 0.20 & 0.04 & 0.03 & 0.02 & 0.00 & 0.00 & 0.09 & $-$0.01 & 0.01 \\
Endoscope action & model-free & 0.24 & 0.05 & 0.01 & 0.01 & 0.00 & 0.00 & 0.01 & $-$0.00 & 0.00 \\
\midrule
Monolithic latent + camera & zero-shot & 0.28 & 0.05 & 0.02 & 0.05 & 0.01 & 0.00 & 0.10 & $-$0.01 & 0.02 \\
Monolithic VQ + camera (SurgWM-style) & zero-shot & 0.12 & 0.02 & 0.02 & 0.02 & 0.00 & 0.00 & 0.00 & 0.00 & 0.00 \\
Slots + camera (FLAM-style) & zero-shot & 0.43 & 0.08 & 0.01 & 0.06 & 0.01 & 0.00 & 0.15 & $-$0.01 & 0.01 \\
EG-SWM, no camera action & zero-shot & 0.30 & 0.02 & 0.02 & 0.01 & 0.00 & $-$0.00 & 0.04 & $-$0.00 & $-$0.00 \\
EG-SWM & zero-shot & 0.29 & 0.01 & 0.02 & 0.01 & 0.00 & $-$0.01 & 0.05 & 0.00 & 0.00 \\
EG-SWM + scene action & zero-shot & 0.33 & 0.01 & 0.02 & 0.01 & 0.00 & $-$0.00 & 0.04 & $-$0.00 & 0.00 \\
EG-SWM, diffusion loss only & zero-shot & 0.11 & 0.00 & $-$0.00 & 0.00 & 0.00 & 0.00 & 0.07 & $-$0.00 & 0.01 \\
EG-SWM, arm actions not grounded & zero-shot & 0.12 & 0.02 & 0.02 & 0.01 & 0.00 & $-$0.00 & 0.04 & $-$0.00 & 0.00 \\
EG-SWM, wrong arm binding & zero-shot & 0.02 & $-$0.01 & 0.01 & 0.00 & 0.00 & $-$0.00 & $-$0.01 & $-$0.00 & 0.00 \\
EG-SWM, size B & zero-shot & 0.28 & 0.00 & 0.02 & 0.01 & 0.00 & $-$0.01 & 0.05 & 0.00 & 0.00 \\
\bottomrule
\end{tabular}
}
\end{table*}}{}

\IfFileExists{tables/e6_detector.tex}{%
\begin{table}[h]
\centering
\caption{\method{} with 25, 50 and 100\,\% of the instrument detections removed at test time.}
\resizebox{\columnwidth}{!}{\begin{tabular}{lcccccc}
\toprule
Detections kept & PSNR $\uparrow$ & SSIM $\uparrow$ & LPIPS $\downarrow$ & $\Delta$PSNR $\uparrow$ & PSNR change & Cases worse \\
\midrule
100\,\% & 22.00\,{\scriptsize[21.43, 22.62]} & 0.691\,{\scriptsize[0.675, 0.708]} & 0.164\,{\scriptsize[0.153, 0.175]} & 2.78 & -- & -- \\
75\,\% & 21.77\,{\scriptsize[21.20, 22.38]} & 0.685\,{\scriptsize[0.669, 0.702]} & 0.170\,{\scriptsize[0.158, 0.181]} & 2.48 & $-$0.23\,{\scriptsize[$-$0.26, $-$0.20]} & 19/19 \\
50\,\% & 21.54\,{\scriptsize[20.98, 22.14]} & 0.680\,{\scriptsize[0.664, 0.697]} & 0.176\,{\scriptsize[0.164, 0.187]} & 2.15 & $-$0.46\,{\scriptsize[$-$0.51, $-$0.43]} & 19/19 \\
0\,\% & 20.94\,{\scriptsize[20.39, 21.51]} & 0.666\,{\scriptsize[0.650, 0.683]} & 0.189\,{\scriptsize[0.178, 0.200]} & 1.33 & $-$1.07\,{\scriptsize[$-$1.17, $-$0.98]} & 19/19 \\
\bottomrule
\end{tabular}
}
\end{table}}{}

\IfFileExists{tables/e1_seeds.tex}{%
\begin{table*}[t]
\centering
\caption{Run-to-run variation: mean and standard deviation over training seeds.}
\resizebox{\textwidth}{!}{\begin{tabular}{lcccccccc}
\toprule
Model & Seeds & PSNR$\uparrow$ & SSIM$\uparrow$ & LPIPS$\downarrow$ & FVD$\downarrow$ & $\Delta$PSNR$\uparrow$ & Cam. gain$\uparrow$ & Locality$\uparrow$ \\
\midrule
EG-SWM & 2 & 22.00 $\pm$ 0.01 & 0.690 $\pm$ 0.001 & 0.164 $\pm$ 0.000 & 187 $\pm$ 1 & 2.77 $\pm$ 0.00 & 0.68 $\pm$ 0.02 & 2.48 $\pm$ 0.09 \\
\bottomrule
\end{tabular}
}
\end{table*}}{}

\IfFileExists{tables/cost.tex}{%
\begin{table*}[t]
\centering
\caption{Cost: training time, throughput and memory, and the time and memory of one 2.8\,s prediction.}
\resizebox{\textwidth}{!}{\begin{tabular}{lcccccccc}
\toprule
Model & Size & Steps & Training (h) & Windows/s & Training memory (GB) & Prediction (s) & VAE decoding (s) & Inference memory (GB) \\
\midrule
None (unconditional) & S & 60000 & 4.6 & 117 & 26.0 & -- & -- & -- \\
Robot state & S & 60000 & 5.2 & 103 & 29.7 & -- & -- & -- \\
Robot state + camera & S & 60000 & 5.2 & 103 & 29.7 & -- & -- & -- \\
Monolithic latent (Genie) & S & 60000 & 5.6 & 96 & 31.4 & -- & -- & -- \\
Monolithic latent + camera & S & 60000 & 5.6 & 96 & 31.8 & 0.29 & 0.37 & 3.7 \\
Monolithic VQ + camera (SurgWM-style) & S & 60000 & 5.6 & 96 & 31.4 & -- & -- & -- \\
Slots + camera (FLAM-style) & S & 60000 & 5.6 & 96 & 31.5 & -- & -- & -- \\
EG-SWM, no camera action & S & 60000 & 5.6 & 96 & 31.4 & -- & -- & -- \\
EG-SWM & S & 60000 & 5.6 & 96 & 31.4 & 0.29 & 0.37 & 3.7 \\
EG-SWM + scene action & S & 60000 & 5.6 & 96 & 31.4 & -- & -- & -- \\
EG-SWM, arm actions not grounded & S & 60000 & 5.6 & 96 & 31.4 & -- & -- & -- \\
\bottomrule
\end{tabular}
}
\end{table*}}{}

\end{document}
